
\documentclass[11pt,a4paper]{article}
\usepackage[spanish]{babel}
\usepackage[utf8]{inputenc}
\usepackage[T1]{fontenc}
\usepackage{amsmath,amssymb}
\usepackage{siunitx}
\usepackage{geometry}
\geometry{margin=2.5cm}

\begin{document}

\title{Ejercicio 2 -- Examen Julio 2024\\Parte (b): efecto del cambio de bielas en un compresor de doble acción}
\author{}
\date{}
\maketitle

\section*{Enunciado (parte b)}

Por una falla en el sistema de lubricación se rompe una biela. No es posible
conseguir el repuesto original, pero el fabricante ofrece una biela que sólo
difiere de la original en su longitud, ya que es \SI{5}{mm} más corta. Sugiere
cambiar las dos bielas para que los pistones trabajen del mismo modo. Indique
si, para una velocidad fija del compresor, el cambio de las bielas podría
afectar el flujo másico de aire o la potencia consumida por el compresor.
Justifique su respuesta cualitativamente.

\section*{1. Modelo geométrico simplificado del cilindro doble acción}

Consideramos un cilindro de doble acción idealizado con las siguientes
características:

\begin{itemize}
  \item El compresor tiene dos cilindros, pero el análisis se hace sobre \emph{un}
        cilindro de doble acción; el comportamiento global se obtiene
        multiplicando por el número de cilindros.
  \item Carrera geométrica (del centro de la cabeza a la culata) constante:
  \[
     S = \SI{120}{mm}.
  \]
  \item Área del pistón $A$ constante en ambas caras (despreciando el volumen
        ocupado por la varilla).
  \item Por revolución, cada cara del pistón barre el mismo volumen
  \[
     V_s = A\,S.
  \]
\end{itemize}

El volumen muerto (``clearance'') se define como
\[
V_3 = V_c,
\]
y el volumen en el comienzo de la admisión es $V_1$. El volumen barrido por
cada lado en una revolución es
\[
V_s = V_1 - V_3.
\]

La fracción de volumen muerto (clearance fraccional) es
\[
C = \frac{V_c}{V_s} = \frac{V_3}{V_1 - V_3}.
\]

\subsection*{1.1. Efecto de acortar las bielas en un doble efecto}

El fabricante propone bielas \SI{5}{mm} más cortas y se sustituyen las \emph{dos}
bielas, de modo que ambas caras del pistón se vean afectadas de forma
simétrica.

Idealizamos el efecto geométrico de la siguiente manera:

\begin{itemize}
  \item La carrera $S = 2R$ viene determinada por el radio de la manivela $R$ y
        \emph{no cambia} al modificar la longitud de la biela $\ell$. Por lo tanto,
        el volumen barrido por cada cara sigue siendo
        \[
          V_s' = V_s'' = V_s = A S.
        \]
  \item La modificación de la longitud de la biela sí altera la posición del
        pistón en el punto muerto superior (PMS) en cada extremo, cambiando los
        volúmenes muertos de ambas cámaras. Denotamos:
        \[
          V_3' = V_3 + \Delta V,\qquad
          V_3'' = V_3 - \Delta V.
        \]
        Es decir, el clearance de una cara aumenta y el de la otra disminuye,
        pero el \emph{promedio} se mantiene:
        \[
          \frac{V_3' + V_3''}{2} = V_3.
        \]
\end{itemize}

Los clearances fraccionales de cada cara son entonces:
\[
C' = \frac{V_3'}{V_s},\qquad
C'' = \frac{V_3''}{V_s},
\]
y por construcción
\[
\frac{C' + C''}{2}
= \frac{1}{2}\left(\frac{V_3'}{V_s} + \frac{V_3''}{V_s}\right)
= \frac{V_3'+V_3''}{2V_s}
= \frac{V_3}{V_s}
= C.
\]

\section*{2. Eficiencia volumétrica de cada cara y eficiencia global}

Para un compresor reciprocante (simple efecto) con proceso politrópico entre
$P_1$ (aspiración) y $P_2$ (descarga), volumen muerto fraccional $C$ y
exponente $k$, la eficiencia volumétrica ideal se escribe como
\[
\eta_v = 1 + C - C\left(\frac{P_2}{P_1}\right)^{1/k}.
\]

Para el cilindro de doble acción, cada cara tiene su propio clearance y, por lo
tanto, su propia eficiencia volumétrica:

\begin{align*}
\eta_v'  &= 1 + C'  - C' \left(\frac{P_2}{P_1}\right)^{1/k},\\
\eta_v'' &= 1 + C'' - C''\left(\frac{P_2}{P_1}\right)^{1/k}.
\end{align*}

Definimos
\[
\beta := \left(\frac{P_2}{P_1}\right)^{1/k}.
\]

\subsection*{2.1. Definición general de eficiencia volumétrica global}

La eficiencia volumétrica global se define como
\[
\eta_{v,\text{tot}} =
\frac{\dot m_{\text{asp real}}}{\dot m_{\text{asp ideal}}}.
\]

En una revolución:

\begin{itemize}
  \item La cara 1 aspira un volumen efectivo
        \[
        V_{\text{asp}}' = \eta_v' V_s,
        \]
        y por tanto una masa
        \[
        m' = \rho_1\,\eta_v' V_s,
        \]
        donde $\rho_1$ es la densidad a la admisión.
  \item La cara 2 aspira
        \[
        V_{\text{asp}}'' = \eta_v'' V_s,
        \qquad
        m'' = \rho_1\,\eta_v'' V_s.
        \]
\end{itemize}

La masa total aspirada por revolución es
\[
m_{\text{tot}} = m' + m'' =
\rho_1\left(\eta_v' V_s + \eta_v'' V_s\right)
= \rho_1 V_s(\eta_v' + \eta_v'').
\]

El volumen ``ideal'' de referencia para la eficiencia global es el volumen
barrido total por revolución, es decir,
\[
V_{s,\text{tot}} = V_s' + V_s'' = 2 V_s.
\]

Luego,
\[
\eta_{v,\text{tot}} =
\frac{m_{\text{tot}}}{\rho_1 V_{s,\text{tot}}}
= \frac{\rho_1 V_s(\eta_v' + \eta_v'')}{\rho_1 \cdot 2 V_s}
= \frac{\eta_v' + \eta_v''}{2}.
\]

Es decir, para un doble efecto \emph{simétrico} (ambas caras con el mismo
volumen barrido),
\[
\boxed{
\eta_{v,\text{tot}} =
\frac{\eta_v' + \eta_v''}{2}.
}
\]

Este resultado no es una hipótesis adicional: se obtiene directamente de la
definición de eficiencia volumétrica global y del hecho de que cada cara tiene
el mismo $V_s$.

En un doble efecto no simétrico ($V_s' \neq V_s''$) la expresión correcta sería
el promedio ponderado
\[
\eta_{v,\text{tot}} =
\frac{\eta_v' V_s' + \eta_v'' V_s''}{V_s' + V_s''}.
\]

\subsection*{2.2. Cálculo de $\eta_{v,\text{tot}}$ en términos de $C'$ y $C''$}

Sustituyendo las expresiones de $\eta_v'$ y $\eta_v''$:

\begin{align*}
\eta_{v,\text{tot}}
&= \frac{1}{2}\left[\eta_v' + \eta_v''\right] \\
&= \frac{1}{2}\left[
(1 + C' - C'\beta) + (1 + C'' - C''\beta)
\right]\\
&= \frac{1}{2}\big[2 + (C' + C'') - (C' + C'')\beta\big]\\[4pt]
&= 1 + \frac{C' + C''}{2} - \frac{C' + C''}{2}\,\beta.
\end{align*}

Definiendo el clearance medio
\[
\bar C := \frac{C' + C''}{2},
\]
se tiene
\[
\eta_{v,\text{tot}} = 1 + \bar C - \bar C\,\beta.
\]

Pero, como vimos en la Sección 1.1, para el cambio de bielas simétrico:
\[
\bar C = \frac{C' + C''}{2} = C.
\]

Por lo tanto,
\[
\boxed{
\eta_{v,\text{tot}} = 1 + C - C\left(\frac{P_2}{P_1}\right)^{1/k},
}
\]
es decir, la eficiencia volumétrica global del cilindro doble acción queda
\textbf{exactamente igual} que antes de cambiar las bielas.

\section*{3. Consecuencias sobre el flujo másico}

El flujo másico del compresor (considerando ambos lados del pistón) a una
velocidad de rotación $N$ (rpm) es
\[
\dot m = \rho_1\,\eta_{v,\text{tot}}\, V_{s,\text{tot}} \frac{N}{60}
       = \rho_1\,\eta_{v,\text{tot}}\, (2 V_s)\frac{N}{60}.
\]

En el problema de la parte (a) se determinó $\eta_v$ y se tomó
\(\eta_{v,\text{tot}} = \eta_v\) (con $C$ original). Con el cambio de bielas
hemos demostrado que:

\begin{itemize}
  \item El volumen barrido total $V_{s,\text{tot}} = 2V_s$ no cambia, puesto que
        la carrera y el área del pistón son las mismas.
  \item La eficiencia volumétrica global $\eta_{v,\text{tot}}$ permanece igual a
        la original.
  \item La densidad a la aspiración $\rho_1$ y la velocidad $N$ son las mismas.
\end{itemize}

Por consiguiente, el flujo másico total a velocidad fija es invariante:
\[
\boxed{\dot m_{\text{nuevo}} = \dot m_{\text{original}}.}
\]

\section*{4. Consecuencias sobre el trabajo y la potencia}

El trabajo específico de compresión (por unidad de masa) para un proceso
politrópico entre $P_1$ y $P_2$ es
\[
w = \frac{k}{k-1}R T_1
\left[
\left(\frac{P_2}{P_1}\right)^{\frac{k-1}{k}} - 1
\right],
\]
que depende sólo de $P_2/P_1$, $k$ y $T_1$, todos ellos datos que no cambian
con la modificación de las bielas.

\begin{itemize}
  \item El trabajo por kilogramo comprimido $w$ es el mismo.
  \item El flujo másico total $\dot m$ no cambia (Sección 3).
\end{itemize}

La potencia indicada del compresor se obtiene como
\[
P_i = \dot m\,w.
\]

Como ambos factores se mantienen, se concluye que:

\[
\boxed{P_{i,\text{nuevo}} = P_{i,\text{original}}.}
\]

Corregida por las mismas eficiencias mecánica y eléctrica, la potencia al eje
y la potencia eléctrica consumida por el compresor también se mantienen
prácticamente iguales.

\section*{5. Conclusión cualitativa para el examen}

Bajo el modelo ideal adoptado:

\begin{itemize}
  \item El cambio simétrico de las dos bielas por otras ligeramente más cortas
        modifica en forma opuesta los volúmenes muertos de las dos cámaras, de
        manera que el \emph{clearance medio} se mantiene y, con él, la
        eficiencia volumétrica global del cilindro doble acción.
  \item El volumen barrido total por revolución no cambia.
  \item El trabajo específico de compresión no cambia, ya que depende sólo del
        radio de compresión y del exponente politrópico.
\end{itemize}

Por lo tanto, para una velocidad fija del compresor:

\[
\boxed{
\text{El cambio de bielas no afecta de manera apreciable ni el flujo másico de aire ni la potencia consumida.}
}
\]

En una discusión más fina podrían aparecer pequeñas diferencias debidas a
asimetrías reales (volumen de la varilla, pérdidas de carga en válvulas,
desajustes geométricos), pero dentro del modelo ideal del examen, el efecto del
cambio de bielas es despreciable sobre la capacidad de caudal y la potencia.

\end{document}
